\subsection*{力矩与角动量的参考点}

相对于固定点 $O$，质点角动量和外力矩分别为
\[
 \vec L_O=\vec r\times m\vec v,
 \qquad
 \vec M_O=\vec r\times\vec F.
\]
外积方向由右手定则确定，大小分别为 $mvr_\perp$ 和 $Fr_\perp$。平行于 $\vec r$ 的力不产生
关于 $O$ 的力矩；但改变参考点会同时改变力臂和角动量，所以“角动量守恒”必须附带明确参考点。
当参考点固定在惯性系中时
\[
 \frac{\dd\vec L_O}{\dd t}=\vec M_O.
\]
若参考点自身加速，直接使用这条形式会遗漏随参考点运动产生的项。

对质点系求和，成对内力若满足第三定律且沿两质点连线，其内力矩相互抵消，得到
$\dd\vec L_O/\dd t=\vec M_{O,\rm ext}$。因此碰撞时把 $O$ 选在支点或接触点，未知冲力的力臂为零，
短暂过程便可直接使用角动量守恒。守恒只要求关于所选点的总外力矩为零，不要求每个外力都为零。

\subsection*{质心平动与绕质心运动}

把 $\vec r_i=\vec R_C+\vec r_i'$、$\vec v_i=\vec V_C+\vec v_i'$ 代入总角动量，可得
\[
 \vec L_O=\vec R_C\times M\vec V_C+\vec L_C.
\]
第一项描述质心作为质量 $M$ 的质点绕 $O$ 的角动量，第二项描述系统相对质心的内部运动。
刚体绕定轴转动时 $\vec L$ 的轴向分量为 $I\omega$，其中
$I=\int r_\perp^2\dd m$；只有转轴沿主惯量轴时，$\vec L$ 才一定与 $\vec\omega$ 平行。

平行轴定理
\[
 I_O=I_C+Md^2
\]
把过质心轴的转动惯量移到与之平行、距离为 $d$ 的轴。它不能用于方向不同的两根轴。计算复合
物体时，可以先把各部分关于自身质心的转动惯量移到公共轴再求和；挖去孔洞则把缺失部分视为
“负质量”作同样运算。

中心力始终沿 $\vec r$，关于力心的力矩为零，故角动量守恒。平面运动中
\[
 L=mr^2\dot\theta=2m\frac{\dd S}{\dd t},
\]
于是面积速度恒定。这不仅给出行星近日点快、远日点慢的定性结论，也把轨道几何和时间联系起来。
若已知势能，还可把径向运动写成含有效势
$U_{\rm eff}(r)=U(r)+L^2/(2mr^2)$ 的一维问题，第二项正是角动量约束产生的“离心势垒”。

